\section{Техническое задание}
\subsection{Основание для разработки}

Основанием для разработки является задание на выпускную квалификационную работу бакалавра «Онлайн-площадка для сдачи вещей в аренду».

\subsection{Цель и назначение разработки}

Основной задачей выпускной квалификационной работы является разработка и тестирование веб-приложения, предоставляющего возможность арендодателям публиковать объявления о сдаче вещей в аренду, а арендаторам - находить, бронировать и оплачивать такие вещи с обеспечением базовых гарантий безопасности сделки.

Задачами данной разработки являются:
\begin{itemize}
\item изучение и анализ предметной области;
\item проектирование интерфейса веб-приложения;
\item проектирование базы данных;
\item проектирование модуля взаимодействия с пользователем;
\item проектирование API веб-приложения.
\end{itemize}
\subsection{Требования к веб-приложению}

\subsubsection{Функциональные требования к веб-приложению}

На основании описанных бизнес-процессов и классификации площадок был сформирован перечень функциональных требований, представленный в таблице \ref{analiz:table}.

\begin{xltabular}{\textwidth}{|c|X|X|}
 \caption{Функциональные требования к системе \label{analiz:table}}\\ \hline
  № & \centrow Функциональное требование & \centrow Описание \\ \hline
 \endfirsthead
 \continuecaption{Продолжение таблицы \ref{analiz:table}}
 ~ № & \centrow Функциональное требование & \centrow Описание \\ \hline 
 \finishhead
 \leftrow 1 & Регистрация и аутентификация & \leftrow Создание учётной записи через email, вход по логину/паролю, восстановление доступа. \\ \hline 
 \leftrow 2 & Верификация пользователей & \leftrow Вводе сведений из документа, удостоверяющего личность, присвоение статуса «Верифицирован». \\ \hline
    \leftrow 3 & Управление профилем & \leftrow Редактирование персональных данных, смена данных, просмотр истории сделок. \\ \hline
    \leftrow 4 & Размещение объявления & \leftrow Создание карточки товара: название, категория, описание, фотографии, цена. \\ \hline
    \leftrow 5 & Управление объявлениями & \leftrow Редактирование, удаление объявления. \\ \hline
    \leftrow 6 & Фильтрация & \leftrow Фильтры по категории, дате доступности. \\ \hline
    \leftrow 7 & Просмотр карточки товара & \leftrow Отображение фото, описания, характеристик, отзывов, цены. \\ \hline
    \leftrow 8 & Онлайн-бронирование & \leftrow Выбор дат аренды, проверка доступности, подтверждение брони арендодателем. \\ \hline
    \leftrow 9 & Онлайн-оплата & \leftrow Оплата банковской картой через платёжный шлюз, удержание средств до подтверждения получения. \\ \hline
    \leftrow 10 & Чат между пользователями & \leftrow Встроенный мессенджер для обсуждения деталей сделки без раскрытия личных контактов. \\ \hline
    \leftrow 11 & Система отзывов & \leftrow Оценки после завершения сделки, накопленный рейтинг пользователей и объявлений. \\ \hline
    \leftrow 12 & Уведомления & \leftrow Push- и email-уведомления о бронировании, оплате, сообщениях, истечении срока аренды. \\ \hline
    \leftrow 13 & Генератор договоров & \leftrow Автоматическое создание договора при оформлении аренды; возможность продолжить сделку только после подписания обеими сторонами. \\ \hline
    \leftrow 14 & Панель аналитики арендодателя & \leftrow Статистика бронирований, доходов объявлений за период. \\ \hline
    \leftrow 15 & Административная панель & \leftrow Модерация объявлений и пользователей, управление категориями, просмотр транзакций, сводная аналитика. \\ \hline
\end{xltabular}

\subsubsection{Требования пользователя к интерфейсу}

Веб-приложение должно включать в себя:
\begin{itemize}
\item регистрацию;
\item авторизацию;
\item профиль пользователя;
\item окно с статистикой по сделкам;
\item встроенный чат между участниками сделки;
\item модуль бронирования с выбором дат аренды и расчётом стоимости;
\item модуль договоров аренды;
\item модуль создания и редактирования объявления.
\item карточку товара;
\item окно с входящими заявками;
\end{itemize}

Схема окна регистрации пользователя представлена на рисунке \ref{reg:image} и состоит из:
\begin{itemize}
\item кнопка для возвращения на предыдущую страницу (1);
\item поле для ввода имени пользователя (2);
\item поле для ввода электронной почты (3);
\item поле для ввода пароля (4);
\item поле для повторного ввода пароля (5);
\item поле для выбора типа пользователя (6);
\item поля для заполнение реквизитов (ИНН, КПП, зависит от типа пользователя) (7);
\item кнопка для регистрации пользователя (8).
\end{itemize}

\begin{figure}[ht]
\center{\includegraphics[width=1\linewidth]{reg}}
\caption{Окно регистрации пользователя}
\label{reg:image}
\end{figure}

Схема окна авторизации пользователя представлена на рисунке \ref{vhod:image} и состоит из:
\begin{itemize}
\item кнопка для возвращения на предыдущую страницу (1);
\item поле для ввода имени пользователя (2);
\item поле для ввода пароля (3);
\item кнопка для авторизации пользователя (4).

\end{itemize}

\begin{figure}[ht]
\center{\includegraphics[width=1\linewidth]{vhod}}
\caption{Окно авторизации пользователя}
\label{vhod:image}
\end{figure}

Схема окна профиля пользователя представлена на рисунке \ref{profile:image} и состоит из:
\begin{itemize}
\item кнопка для возвращения на предыдущую страницу (1);
\item поле для ввода имени пользователя (2);
\item поле для ввода электронной почты (3);
\item поля для заполнение реквизитов (ИНН, КПП, зависит от типа пользователя) (4);
\item кнопка для сохранения изменений (5);
\item поле для ввода текущего пароля (6);
\item поле для ввода нового пароля (7);
\item поле для повторного ввода нового пароля (8);
\item кнопка для смены пароля (9).
\end{itemize}

\begin{figure}[H]
\center{\includegraphics[width=1\linewidth]{profile}}
\caption{Окно профиля пользователя}
\label{profile:image}
\end{figure}

Схема окна с статистикой по сделкам представлена на рисунке \ref{stat:image} и состоит из:
\begin{itemize}
\item поле для вывода общего количества объявлений созданные пользователем (1);
\item поле для вывода количества бронирований у пользователя (2);
\item поле для вывода количества забронированных товаров пользователем (3);
\item поле для вывода общий прибыли за аренду товаров пользователя (4);
\item поле для вывода объявлений созданных пользователем (5);
\item поле для вывода объявлений находящихся в аренде (6);
\item поле для вывода объявлений с истекшим сроком аренды (7);
\item поле для вывода объявлений забронированных пользователем (8).
\end{itemize}

\begin{figure}[H]
\center{\includegraphics[width=1\linewidth]{stat}}
\caption{Окно с статистикой по сделкам}
\label{stat:image}
\end{figure}

Схема окна встроенного чат между участниками сделки представлена на рисунке \ref{chat:image} и состоит из:
\begin{itemize}
\item кнопка для возвращения на предыдущую страницу (1);
\item кнопка для отправки сообщения (2);
\item поле для ввода сообщения (3);
\item кнопка для загрузки изображения в сообщение (4).
\end{itemize}

\begin{figure}[H]
\center{\includegraphics[width=1\linewidth]{chat}}
\caption{Окно встроенного чат между участниками сделки}
\label{chat:image}
\end{figure}

Схема интерфейса модуля бронирования с выбором дат аренды и расчётом стоимости представлена на рисунке  и состоит из:
\begin{itemize}
\item поле для ввода даты начала аренды (1);
\item поле для ввода даты окончания аренды (2);
\item поле для вывода статуса брони (3);
\item кнопка для подписи договора аренды через демонстрационный механизм ЭЦП (4);
\item кнопка для отправки заявки на бронь (5).
\end{itemize}

\begin{figure}[H]
\center{\includegraphics[width=1\linewidth]{booking}}
\caption{Окно встроенного чат между участниками сделки}
\label{booking:image}
\end{figure}

Схема интерфейса модуля модуля договоров аренды представлена на рисунке \ref{contract:image} и состоит из:
\begin{itemize}
\item поле для вывода данных договора(1);
\item поле для вывода сводки по договору (2);
\item кнопка для установки PDF-файла договора аренды (3);
\item поле для ввода ФИО подписанта(4);
\item поле для ввода pin сертификата (5);
\item кнопка для подписывание договора (6).
\end{itemize}

\begin{figure}[H]
\center{\includegraphics[width=1\linewidth]{contract}}
\caption{Интерфейс модуля договоров аренды}
\label{contract:image}
\end{figure}

Схема интерфейса модуля создания и редактирования объявления представлена на рисунке \ref{item:image} и состоит из:
\begin{itemize}
\item кнопка для возвращения на предыдущую страницу (1);
\item поле для ввода наименования товара (2);
\item поле для ввода описания товара (3);
\item поле для ввода цены товара (4);
\item поле для выбора категории товара (5);
\item поле для загрузки изображения (6);
\item кнопка для публикации объявления (7).
\end{itemize}

\begin{figure}[H]
\center{\includegraphics[width=1\linewidth]{item}}
\caption{Интерфейс модуля создания и редактирования объявления}
\label{item:image}
\end{figure}

Схема окна карточки товара представлена на рисунке \ref{item_card:image} и состоит из:
\begin{itemize}
\item поле для вывода изображений (1);
\item кнопка для отправки заявки на аренду товара (2);
\item кнопка для перенаправления пользователя в чат с владельцем объявления (3);
\item поле для вывода отзывов по объявлению (4).
\end{itemize}

\begin{figure}[H]
\center{\includegraphics[width=1\linewidth]{item_card}}
\caption{Окно карточки товара}
\label{item_card:image}
\end{figure}

Схема окна с входящими заявками представлена на рисунке \ref{order:image} и состоит из:
\begin{itemize}
\item кнопка для подтверждения заявки (1);
\item кнопка для отклонения заявки (2);
\item поле для вывода входящих заявок (3);
\item поле для вывода заявок оставленные пользователем (4);
\item кнопка для открытия объявления (5).
\end{itemize}

\begin{figure}[H]
\center{\includegraphics[width=1\linewidth]{order}}
\caption{Окно с входящими заявками}
\label{order:image}
\end{figure}

\subsection{Моделирование вариантов использования}

Для разрабатываемого веб-приложения была реализована модель, обеспечивающая наглядное представление вариантов использования программы с помощью унифицированного языка визуального моделирования UML. На основании анализа предметной области в программе должны быть реализованы следующие прецеденты:
\begin{enumerate}
    \item Регистрация 
    \item Авторизация 
    \item Изменение пароля пользователя. 
    \item Просмотр профиля пользователя. 
    \item Редактирование профиля пользователя. 
    \item Просмотр каталога вещей. 
    \item Фильтрация объявлений по категории и дате. 
    \item Создание собственных объявлений. 
    \item Редактирование собственных объявлений. 
    \item Создание бронирования вещи с автоматическим расчетом стоимости.
    \item Удаление собственных объявлений.
    \item Загрузка изображений в объявление. 
    \item Удаление изображений в объявлении. 
    \item Просмотр карточки вещи,ее рейтинга,отзывов и занятых периодов аренды. 
    \item Управление статусами бронирования:подтверждение,отмена. 
    \item Просмотр входящих и исходящих бронирований пользователя. 
    \item Оплата бронирования через демонстрационный механизм СБП. 
    \item Формирование и подписание договора аренды участниками сделки.
    \item Создание отзывов по завершенным или оплаченным бронированиям.
    \item Чат между пользователями с возможностью отправить изображение. 
    \item Помощь на часто задаваемые вопросы.
\end{enumerate}

На рисунке \ref{diaprec:image} изображены прецеденты для пользователя.
\begin{figure}[H]
\center{\includegraphics[width=1\linewidth]{diaprec}}
\caption{Диаграмма прецедентов для пользователя}
\label{diaprec:image}
\end{figure}

\subsubsection{Регистрация}
Заинтересованные лица и их требования: пользователь желает создать учетную запись для получения доступа к функциям сервиса аренды вещей.
Предусловие: пользователь находится на странице регистрации.
Постусловие: в системе создана новая учетная запись пользователя.
\begin{enumerate}
    \item Пользователь открывает форму регистрации.
    \item Пользователь вводит имя пользователя, email, пароль и выбирает тип аккаунта.
    \item Веб-приложение передает введенные данные на сервер.
    \item Сервер проверяет корректность данных, уникальность email и длину пароля.
    \item Сервер создает учетную запись пользователя и сохраняет ее в базе данных.
    \item Веб-приложение сообщает пользователю об успешной регистрации.
\end{enumerate}
 
\subsubsection{Авторизация}
Заинтересованные лица и их требования: пользователь желает войти в систему для доступа к личному кабинету, объявлениям, бронированиям и сообщениям.
Предусловие: пользователь зарегистрирован в системе и находится на странице входа.
Постусловие: пользователь авторизован и получает доступ к функциям личного кабинета.
Основной успешный сценарий:
\begin{enumerate}
    \item Пользователь открывает форму авторизации.
    \item Пользователь вводит email или имя пользователя и пароль.
    \item Приложение передает данные авторизации на сервер.
    \item Сервер проверяет соответствие введенных данных учетной записи.
    \item Сервер подтверждает авторизацию пользователя.
    \item Веб-приложение открывает пользователю доступ к личному кабинету.
\end{enumerate}

\subsubsection{Изменение пароля пользователя}
Заинтересованные лица и их требования: пользователь желает изменить пароль для защиты своей учетной записи.

Предусловие: пользователь авторизован в системе.

Постусловие: пароль пользователя изменен, дальнейшая авторизация выполняется с новым паролем.

Основной успешный сценарий:
\begin{enumerate}
    \item Пользователь восстанавливает пароль без входа в систему, он открывает форму «Забыли пароль?».
    \item Пользователь вводит email, указанный при регистрации, и нажимает кнопку отправки кода.
    \item Приложение передает email на сервер.
    \item Сервер находит пользователя по email, генерирует шестизначный код подтверждения и сохраняет его во временном хранилище.
    \item Сервер отправляет код подтверждения на email пользователя.
    \item Пользователь вводит полученный код и новый пароль.
    \item Приложение передает email, код подтверждения и новый пароль на сервер.
    \item Сервер проверяет корректность кода и срок его действия.
    \item Сервер проверяет, что новый пароль соответствует требованиям, в частности содержит не менее 6 символов.
    \item Сервер сохраняет новый пароль пользователя и удаляет использованный код подтверждения.
    \item Приложение сбрасывает текущую сессию пользователя и предлагает войти заново с новым паролем.
\end{enumerate}
\subsubsection{Просмотр профиля пользователя}
Заинтересованные лица и их требования: пользователь желает просмотреть данные своего аккаунта и проверить заполненность профиля.

Предусловие: пользователь авторизован в системе.

Постусловие: пользователю отображаются данные его профиля.

Основной успешный сценарий:

\begin{enumerate}
    \item Пользователь открывает раздел «Профиль».
    \item Приложение отправляет запрос на сервер для получения данных текущего пользователя.
    \item Сервер находит профиль пользователя в базе данных.
    \item Сервер возвращает данные профиля: имя пользователя, email, тип аккаунта и реквизиты.
    \item Приложение отображает полученные данные пользователю.
\end{enumerate}

\subsubsection{Редактирование профиля пользователя}
Заинтересованные лица и их требования: пользователь желает изменить данные профиля с учетом типа аккаунта.

Предусловие: пользователь авторизован и открыл раздел «Профиль».

Постусловие: данные профиля обновлены, признак заполненности профиля пересчитан.

Основной успешный сценарий:
\begin{enumerate}
    \item Пользователь нажимает кнопку редактирования профиля.
    \item Пользователь изменяет необходимые поля.
    \item Для физического лица указываются ФИО, серия и номер паспорта, ИНН.
    \item Приложение передает обновленные данные на сервер.
    \item Сервер проверяет корректность и заполненность обязательных полей.
    \item Сервер сохраняет изменения в базе данных.
    \item Приложение отображает обновленный профиль пользователя.
\end{enumerate}
\subsubsection{Просмотр каталога вещей}
Заинтересованные лица и их требования: пользователь желает просмотреть список вещей, доступных для аренды.

Предусловие: пользователь открыл страницу каталога вещей.

Постусловие: на странице отображается список объявлений о вещах.

Основной успешный сценарий:
\begin{enumerate}
    \item Пользователь переходит в раздел каталога вещей.
    \item Приложение отправляет запрос на сервер для получения списка объявлений.
    \item Сервер выбирает из базы данных доступные вещи.
    \item Сервер возвращает сведения о вещах: название, описание, цену, категорию, рейтинг и изображения.
    \item Приложение отображает полученный список объявлений.
    \item Пользователь просматривает каталог и выбирает интересующую вещь.
\end{enumerate}
\subsubsection{Фильтрация объявлений по категории и дате}
Заинтересованные лица и их требования: пользователь желает отобрать объявления по нужной категории и отсортировать их по дате создания.

Предусловие: пользователь находится в каталоге вещей.

Постусловие: каталог обновлен и содержит объявления, соответствующие выбранным параметрам.

Основной успешный сценарий:
\begin{enumerate}
    \item Пользователь выбирает категорию вещи в интерфейсе каталога.
    \item Пользователь выбирает сортировку по дате создания.
    \item Приложение передает выбранные параметры фильтрации и сортировки на сервер.
    \item Сервер формирует запрос к базе данных с учетом категории.
    \item Сервер сортирует найденные объявления по дате создания.
    \item Сервер передает результат в приложение.
    \item Приложение обновляет список отображаемых объявлений.
\end{enumerate}
\subsubsection{Создание собственных объявлений}
Заинтересованные лица и их требования: пользователь желает разместить собственную вещь для аренды.

Предусловие: пользователь авторизован, его профиль заполнен.

Постусловие: новое объявление создано и доступно в системе.

Основной успешный сценарий:
\begin{enumerate}
    \item Пользователь открывает форму создания объявления.
    \item Пользователь вводит название, описание, цену за день и выбирает категорию.
    \item Пользователь указывает статус вещи.
    \item Приложение передает данные объявления на сервер.
    \item Сервер проверяет авторизацию пользователя и заполненность профиля.
    \item Сервер сохраняет объявление и назначает пользователя владельцем вещи.
    \item Приложение отображает созданное объявление в списке вещей.
\end{enumerate}
\subsubsection{Редактирование собственных объявлений}
Заинтересованные лица и их требования: владелец вещи желает изменить данные ранее созданного объявления.

Предусловие: пользователь авторизован и является владельцем объявления.

Постусловие: данные объявления обновлены.

Основной успешный сценарий:
\begin{enumerate}
    \item Пользователь открывает список своих объявлений.
    \item Пользователь выбирает объявление и нажимает кнопку редактирования.
    \item Пользователь изменяет название, описание, цену, категорию или статус вещи.
    \item Приложение передает обновленные данные на сервер.
    \item Сервер проверяет, что текущий пользователь является владельцем объявления.
    \item Сервер сохраняет изменения в базе данных.
    \item Приложение отображает обновленную карточку объявления.
\end{enumerate}
\subsubsection{Создание бронирования вещи с автоматическим расчетом стоимости}
Заинтересованные лица и их требования: пользователь желает забронировать вещь и сразу увидеть итоговую стоимость аренды.

Предусловие: пользователь авторизован, профиль заполнен, открыта карточка чужой вещи.

Постусловие: бронирование создано, итоговая стоимость рассчитана автоматически.

Основной успешный сценарий:
\begin{enumerate}
    \item Пользователь нажимает кнопку бронирования в карточке вещи.
    \item Пользователь выбирает дату начала и дату окончания аренды.
    \item Приложение передает данные бронирования на сервер.
    \item Сервер проверяет, что пользователь не является владельцем вещи.
    \item Сервер проверяет заполненность профиля пользователя.
    \item Сервер рассчитывает количество дней аренды.
    \item Сервер умножает количество дней на цену вещи за день.
    \item Сервер сохраняет бронирование со статусом рассматриваемое.
    \item Приложение отображает созданную заявку и рассчитанную стоимость.
\end{enumerate}
\subsubsection{Удаление собственных объявлений}
Заинтересованные лица и их требования: владелец вещи желает удалить свое объявление из системы.

Предусловие: пользователь авторизован и является владельцем объявления.

Постусловие: объявление удалено из базы данных и больше не отображается в каталоге.

Основной успешный сценарий:
\begin{enumerate}
    \item Пользователь открывает список своих объявлений.
    \item Пользователь выбирает объявление для удаления.
    \item Пользователь подтверждает удаление.
    \item Приложение отправляет запрос на сервер.
    \item Сервер проверяет, что пользователь является владельцем объявления.
    \item Сервер удаляет объявление из базы данных.
    \item Приложение обновляет список объявлений пользователя.
\end{enumerate}
\subsubsection{Загрузка изображений в объявление}
Заинтересованные лица и их требования: владелец объявления желает добавить изображения вещи для более подробного представления объявления.

Предусловие: пользователь авторизован и является владельцем объявления.

Постусловие: изображение загружено и связано с объявлением.

Основной успешный сценарий:
\begin{enumerate}
    \item Пользователь открывает свое объявление.
    \item Пользователь выбирает изображение для загрузки.
    \item Приложение передает файл изображения и идентификатор вещи на сервер.
    \item Сервер проверяет, что вещь принадлежит текущему пользователю.
    \item Сервер сохраняет изображение в файловом хранилище.
    \item Сервер связывает изображение с объявлением.
    \item Приложение отображает загруженное изображение в карточке вещи.
\end{enumerate}
\subsubsection{Удаление изображений в объявлении}
Заинтересованные лица и их требования: владелец объявления желает удалить лишнее или устаревшее изображение вещи.

Предусловие: пользователь авторизован и является владельцем объявления, к объявлению добавлено изображение.

Постусловие: выбранное изображение удалено из объявления.

Основной успешный сценарий:
\begin{enumerate}
    \item Пользователь открывает свое объявление.
    \item Пользователь выбирает изображение для удаления.
    \item Пользователь подтверждает удаление изображения.
    \item Приложение отправляет запрос на сервер.
    \item Сервер проверяет, что изображение относится к вещи текущего пользователя.
    \item Сервер удаляет изображение.
    \item Приложение обновляет список изображений объявления.
\end{enumerate}
\subsubsection{Просмотр карточки вещи, ее рейтинга, отзывов и занятых периодов аренды}
Заинтересованные лица и их требования: пользователь желает получить подробную информацию о вещи перед арендой.

Предусловие: пользователь находится в каталоге вещей.

Постусловие: открыта карточка вещи с подробной информацией.

Основной успешный сценарий:
\begin{enumerate}
    \item Пользователь выбирает вещь в каталоге.
    \item Приложение отправляет запрос на сервер по идентификатору вещи.
    \item Сервер получает данные вещи, владельца, категории и изображений.
    \item Сервер получает рейтинг вещи и количество отзывов.
    \item Сервер получает список отзывов о вещи.
    \item Сервер определяет занятые периоды аренды по активным бронированиям.
    \item Приложение отображает карточку вещи, рейтинг, отзывы и занятые даты.
\end{enumerate}

\subsubsection{Управление статусами бронирования: подтверждение, отмена}
Заинтересованные лица и их требования: владелец и арендатор желают управлять заявкой на аренду.

Предусловие: бронирование уже создано.

Постусловие: статус бронирования изменен на принято или отменено.

Основной успешный сценарий:
\begin{enumerate}
    \item Владелец вещи открывает входящее бронирование.
    \item Владелец нажимает кнопку подтверждения.
    \item Приложение отправляет запрос на сервер.
    \item Сервер проверяет, что текущий пользователь является владельцем вещи.
    \item Сервер изменяет статус бронирования на принято.
    \item Сервер отправляет арендатору уведомление о подтверждении.
    \item При отмене участник сделки открывает бронирование.
    \item Пользователь нажимает кнопку отмены.
    \item Сервер проверяет, что пользователь является арендатором или владельцем.
    \item Сервер изменяет статус бронирования на отменено.
\end{enumerate}
\subsubsection{Просмотр входящих и исходящих бронирований пользователя}
Заинтересованные лица и их требования: пользователь желает просматривать заявки, созданные им самим, и заявки, поступившие на его вещи.

Предусловие: пользователь авторизован в системе.

Постусловие: пользователю отображается список входящих и исходящих бронирований.

Основной успешный сценарий:
\begin{enumerate}
    \item Пользователь открывает раздел бронирований.
    \item Приложение отправляет запрос на сервер.
    \item Сервер выбирает бронирования, где пользователь является арендатором или владельцем вещи.
    \item При выборе входящих бронирований сервер возвращает заявки на вещи пользователя.
    \item При выборе исходящих бронирований сервер возвращает заявки, созданные пользователем.
    \item Приложение отображает список бронирований.
    \item Пользователь видит вещь, даты аренды, стоимость, статус бронирования и статус оплаты.
\end{enumerate}
\subsubsection{Оплата бронирования через демонстрационный механизм СБП}
Заинтересованные лица и их требования: арендатор желает оплатить подтвержденное бронирование через демонстрационный механизм СБП.

Предусловие: бронирование подтверждено владельцем и имеет статус confirmed.

Постусловие: бронирование получает статус оплаты paid.

Основной успешный сценарий:
\begin{enumerate}
    \item Арендатор открывает подтвержденное бронирование.
    \item Пользователь нажимает кнопку оплаты.
    \item Приложение отправляет запрос на сервер для запуска оплаты.
    \item Сервер проверяет, что бронирование подтверждено и еще не оплачено.
    \item Сервер создает демонстрационную СБП-сессию.
    \item Сервер возвращает сумму, получателя, банк, ссылку и QR-данные.
    \item Пользователь выполняет демонстрационную оплату.
    \item Пользователь подтверждает оплату на сайте.
    \item Сервер изменяет статус оплаты на оплачено.
    \item Владелец вещи получает уведомление об оплате.
\end{enumerate}
\subsubsection{Формирование и подписание договора аренды участниками сделки}
Заинтересованные лица и их требования: арендатор и владелец желают сформировать и подписать договор аренды по бронированию.

Предусловие: бронирование подтверждено или завершено.

Постусловие: договор сформирован, после подписи обеих сторон отмечен как подписанный.

Основной успешный сценарий:
\begin{enumerate}
    \item Участник сделки открывает договор по бронированию.
    \item Приложение отправляет запрос на сервер.
    \item Сервер проверяет, что пользователь является участником сделки.
    \item Сервер проверяет статус бронирования.
    \item Сервер создает договор, если он еще не был создан.
    \item Сервер формирует номер договора и файл с данными аренды.
    \item Пользователь вводит имя подписанта и PIN демонстрационной подписи.
    \item Сервер сохраняет подпись участника.
    \item Второй участник выполняет подписание аналогичным образом.
    \item После подписания обеими сторонами сервер отмечает договор как подписанный.
\end{enumerate}
\subsubsection{Создание отзывов по завершенным или оплаченным бронированиям}
Заинтересованные лица и их требования: арендатор и владелец желают обсудить детали аренды во встроенном чате.

Предусловие: пользователь авторизован в системе.

Постусловие: сообщение отправлено, сохранено и доступно участникам чата.

Основной успешный сценарий:
\begin{enumerate}
    \item Пользователь открывает карточку вещи или раздел чатов.
    \item Пользователь выбирает собеседника и начинает диалог.
    \item Приложение передает на сервер идентификатор собеседника и при необходимости идентификатор вещи.
    \item Сервер проверяет наличие существующего чата между пользователями.
    \item Если чат отсутствует, сервер создает новый чат.
    \item Пользователь вводит текст сообщения или прикрепляет изображение.
    \item Приложение отправляет сообщение на сервер.
    \item Сервер сохраняет сообщение в базе данных.
    \item Сервер отправляет уведомление получателю.
    \item Приложение отображает новое сообщение в чате.
\end{enumerate}

\subsubsection{Чат между пользователями с возможностью отправить изображение}

Заинтересованные лица и их требования: арендатор и владелец желают обсудить детали аренды во встроенном чате.

Предусловие: пользователь авторизован в системе.

Постусловие: сообщение отправлено, сохранено и доступно участникам чата.

Основной успешный сценарий:
\begin{enumerate}
    \item Пользователь открывает карточку вещи или раздел чатов.
    \item Пользователь выбирает собеседника и начинает диалог.
    \item Приложение передает на сервер идентификатор собеседника и при необходимости идентификатор вещи.
    \item Сервер проверяет наличие существующего чата между пользователями.
    \item Если чат отсутствует, сервер создает новый чат.
    \item Пользователь вводит текст сообщения или прикрепляет изображение.
    \item Сервер сохраняет сообщение в базе данных.
    \item Сервер отправляет уведомление получателю.
    \item Приложение отображает новое сообщение в чате.
\end{enumerate}

\subsubsection{Помощь на часто задаваемые вопросы}

Заинтересованные лица и их требования: пользователь желает получить справочную информацию по основным действиям в системе аренды вещей: созданию объявления, заполнению профиля, оформлению аренды, оплате, переписке и просмотру чатов.

Предусловие: пользователь находится на сайте и открыл раздел помощи или окно поддержки.

Постусловие: пользователь получает список часто задаваемых вопросов и ответ с пошаговой инструкцией по выбранной теме.

Основной успешный сценарий:

\begin{enumerate}
    \item Пользователь открывает раздел помощи или окно поддержки на сайте.
    \item Приложение отправляет запрос на сервер для получения списка часто задаваемых вопросов.
    \item Сервер получает подготовленный список FAQ-вопросов из системы.
    \item Сервер передает в приложение список вопросов и связанных с ними ответов.
    \item Приложение отображает пользователю доступные темы помощи.
    \item Пользователь выбирает интересующий вопрос, например о создании объявления, заполнении профиля, бронировании, оплате или чатах.
    \item Приложение отображает ответ по выбранной теме.
    \item Пользователь просматривает пошаговую инструкцию.
    \item При необходимости пользователь переходит в соответствующий раздел сайта и выполняет действие по инструкции.
\end{enumerate}

\subsection{Требования к оформлению документации}

Разработка программной документации и программного изделия должна производиться согласно ГОСТ 19.102-77 и ГОСТ 34.601-90. Единая система программной документации.